\documentclass[11pt]{article}
\usepackage[margin=1in]{geometry}
\usepackage{amsmath,amssymb}
\usepackage{graphicx}
\usepackage{booktabs}
\usepackage{hyperref}
\usepackage{siunitx}
\title{Replication Report:\\ \emph{Many-Body Order Parameters for Multipoles in Solids}\\
\normalsize Kang, Shiozaki, Cho --- arXiv:1812.06999 (Phys.\ Rev.\ B \textbf{100}, 245134, 2019)}
\author{Automated replication (OpenClaw subagent)}
\date{2026-07-19}
\begin{document}
\maketitle

\section{Summary}
We replicate the central numerical claims of Kang, Shiozaki, and Cho, who
introduce gauge-invariant, real-space \emph{many-body} order parameters for the
electric multipoles of a crystalline insulator and show they equal the physical
corner charge. Working entirely with a minimal tight-binding model and linear
algebra (no DFT, no paid services), we reproduce the quantized quadrupole moment
($0$ vs $1/2$), the sharp phase transition at the Wannier-gap closing, the
vanishing of the dipole in the quadrupole phase, and the Thouless-pump behavior.
\textbf{All five machine-checkable claims pass.}

\section{Theory reproduced}
The paper generalizes Resta's many-body polarization to multipoles. For a
ground state $|\Psi\rangle$ on a torus with periodic boundary conditions,
\begin{align}
P_x &= \tfrac{1}{2\pi}\,\mathrm{Im}\ln\langle \hat U_1\rangle, &
 \hat U_1 &= \exp\!\Big(2\pi i \sum_x \tfrac{x}{L_x}\,\hat n(x)\Big), \tag{Eq.~1}\\
Q_{xy} &= \tfrac{1}{2\pi}\,\mathrm{Im}\ln\langle \hat U_2\rangle, &
 \hat U_2 &= \exp\!\Big(2\pi i \sum_{\mathbf r} \tfrac{xy}{L_xL_y}\,\hat n(\mathbf r)\Big). \tag{Eq.~2}
\end{align}
The key claim is $\langle\hat U_a\rangle=\exp(2\pi i\,q_c)$, i.e.\ the phase
equals the corner charge / physical multipole, even when quantizing symmetries
are broken, in contrast to the nested-Wilson-loop index which only works under
symmetry.

\paragraph{Model (Eq.~8 / D5).} The non-interacting Benalcazar--Bernevig--Hughes
quadrupole insulator, four orbitals per cell on a square lattice with $\pi$-flux:
\begin{equation}
h(\mathbf k)=\big(\gamma_x+\lambda_x\cos k_x\big)\Gamma_4+\lambda_x\sin k_x\,\Gamma_3
+\big(\gamma_y+\lambda_y\cos k_y\big)\Gamma_2+\lambda_y\sin k_y\,\Gamma_1+\delta\,\Gamma_0,
\end{equation}
with $\Gamma_0=\tau_3\!\otimes\!\tau_0$, $\Gamma_i=-\tau_2\!\otimes\!\tau_i$
($i=1,2,3$), $\Gamma_4=\tau_1\!\otimes\!\tau_0$. At half filling and $\delta=0$
the ground state is a topological quadrupole insulator ($q_{xy}=1/2$) for
$|\gamma_x/\lambda_x|<1$ and $|\gamma_y/\lambda_y|<1$; the physical gap closes
only when both ratios equal $1$, while the Wannier gap (hence the $Q_{xy}$
transition) closes when \emph{either} ratio equals $1$.

\paragraph{Evaluation.} For a Slater determinant we use the determinant identity
\begin{equation}
\big\langle\Psi\big|\exp\!\big(i\textstyle\sum_{\mathbf r}\phi(\mathbf r)\hat n_{\mathbf r}\big)\big|\Psi\big\rangle
=\det\!\big(P^{\dagger}\,\mathrm{diag}(e^{i\phi(\mathbf r)})\,P\big),
\end{equation}
with $P$ the matrix of occupied single-particle eigenvectors on an $L\times L$
torus, obtained by inverse Fourier transform of the Bloch eigenstates.
$Q_{xy}=\tfrac{1}{2\pi}\mathrm{Im}\ln\det(\cdot)$, referenced to the atomic
trivial limit ($\lambda\to0$) to fix the coordinate origin (the paper proves
invariance under adding a trivial atomic band).

\section{Results}

\begin{table}[h]\centering
\caption{C1/C2 --- quantized quadrupole ($\delta=0$, $\lambda=1$). Trivial
$\gamma=1.5$; topological $\gamma=0.5$. $Q_{xy}$ referenced to the atomic limit.}
\begin{tabular}{ccccc}
\toprule
$L$ & $Q_{xy}$ (trivial) & $Q_{xy}$ (topo) & $|\langle\hat U_2\rangle|$ triv & $|\langle\hat U_2\rangle|$ topo\\
\midrule
6  & $0.0000$ & $0.5000$ & \num{1.78e-1} & \num{7.9e-4}\\
8  & $0.0000$ & $0.5000$ & \num{1.10e-1} & \num{2.4e-4}\\
10 & $0.0000$ & $0.5000$ & \num{7.21e-2} & \num{7.2e-5}\\
12 & $0.0000$ & $0.5000$ & \num{4.91e-2} & \num{2.2e-5}\\
\bottomrule
\end{tabular}
\end{table}

\begin{table}[h]\centering
\caption{Quantitative comparison to the paper.}
\begin{tabular}{lccc}
\toprule
Quantity & Paper & This work & Match\\
\midrule
$Q_{xy}$ trivial & $0$ & $0.0000$ & exact\\
$Q_{xy}$ topological & $1/2$ & $0.5000$ & exact\\
Transition location & $|\gamma/\lambda|=1$ & $\gamma_y=1.00$ & exact\\
$|\langle\hat U_2\rangle|$ at boundary & $\to 0$ & $\sim\!10^{-4}$ dip & qual.\\
Pump quantized pts ($\delta=0$) & $0,\ 1/2$ & $0.0000,\ 0.5000$ & exact\\
Dipole in quad.\ phase & $0$ & $0.0000$ & exact\\
\bottomrule
\end{tabular}
\end{table}

\paragraph{C3 --- sharp transition.} Sweeping $\gamma_y$ at $\gamma_x=0.5$,
$\lambda=1$, $\delta=0$ ($L=10$), $|Q_{xy}|=0.5$ for $\gamma_y<1$ and $0$ for
$\gamma_y>1$, jumping exactly at $\gamma_y=1$ (Wannier-gap closing), while the
physical energy gap stays open --- matching Fig.~1(c).

\paragraph{C4 --- polarization.} In both phases the many-body dipole
$P_x=P_y=0.0000$, satisfying the well-definedness condition for $Q_{xy}$.

\paragraph{C5 --- Thouless pump (Eq.~9).} Along the isotropic pump, $Q_{xy}$
pins to $0.5$ at $\theta=\pi/2$ (topological, $\delta=0$) and $0$ at
$\theta=3\pi/2$ (trivial, $\delta=0$), varying continuously in between as
$\delta\neq0$ de-quantizes the moment --- consistent with Fig.~2(a), where the
phase of $\langle\hat U_2\rangle$ tracks the corner charge.

\begin{figure}[h]\centering
\includegraphics[width=\textwidth]{../work/replication_figures.png}
\caption{Left: C3 transition (mirror of Fig.~1c). Right: C5 isotropic Thouless
pump (mirror of Fig.~2a); dotted lines mark the $\delta=0$ quantized crossings.}
\end{figure}

\section{Verdict}
\textbf{Reproduced.} The free-fermion core of the paper --- the quantized
quadrupole, its sharp Wannier-gap transition, the vanishing dipole, and the
Thouless-pump quantization --- is faithfully recovered by a minimal
tight-binding + determinant calculation.
\begin{itemize}
\item \textbf{Coverage: 6/10.} The main-text non-interacting results are fully
covered. Not covered (out of scope, minimal model): interacting/bosonic ground
states (the paper's headline generality), the anomalous quadrupole insulator
(Fig.~1d), the partial-region operators $V(l)$ (Fig.~3), the octupole, and the
explicit nested-Wilson-loop cross-check.
\item \textbf{Agreement: 9/10.} Every reproduced quantity matches the paper
exactly ($0$, $1/2$, transition at $|\gamma/\lambda|=1$). One point withheld
because absolute $Q_{xy}$ required an atomic-limit reference to fix the origin
(gauge/coordinate convention), and $|\langle\hat U_2\rangle|$ scaling was only
checked qualitatively.
\end{itemize}

\section{Reproducibility}
\texttt{python3 code/run\_all.py work} then \texttt{python3 code/make\_figures.py work}.
See \texttt{workflow.md}; open questions in \texttt{open\_questions.json};
issues in \texttt{failure\_analysis.md}.
\end{document}
